% ============================================================
% fig_st_problem.tex  ST边界叠加问题直观示意
% 用途：研究背景幻灯片（Slide 4）右侧说明图
% 用法：% ============================================================
% fig_st_problem.tex  ST边界叠加问题直观示意
% 用途：研究背景幻灯片（Slide 4）右侧说明图
% 用法：% ============================================================
% fig_st_problem.tex  ST边界叠加问题直观示意
% 用途：研究背景幻灯片（Slide 4）右侧说明图
% 用法：% ============================================================
% fig_st_problem.tex  ST边界叠加问题直观示意
% 用途：研究背景幻灯片（Slide 4）右侧说明图
% 用法：\input{figures/fig_st_problem.tex}
% ============================================================
\begin{tikzpicture}[scale=0.9]
  % 坐标轴
  \draw[->, thick] (0,0) -- (4.3,0) node[right]{\footnotesize $t$/s};
  \draw[->, thick] (0,0) -- (0,3.3) node[above]{\footnotesize $s$/m};
  % 障碍物 1（实心红）
  \fill[red!50] (0.5,0.35) rectangle (1.7,1.05);
  \node[font=\tiny, white] at (1.1,0.7) {\textbf{障碍1}};
  % 障碍物1 安全缓冲（虚线）
  \draw[red!70, dashed, thick] (0.2,0.05) rectangle (2.0,1.35);
  \node[font=\tiny, red!70!black] at (1.7,0.15) {缓冲1};
  % 障碍物 2（实心蓝）
  \fill[blue!50] (1.3,1.15) rectangle (2.5,1.95);
  \node[font=\tiny, white] at (1.9,1.55) {\textbf{障碍2}};
  % 障碍物2 安全缓冲（虚线）
  \draw[blue!70, dashed, thick] (1.0,0.85) rectangle (2.8,2.25);
  \node[font=\tiny, blue!70!black] at (2.5,2.35) {缓冲2};
  % 叠加区高亮（紫）
  \fill[purple!40, opacity=0.75] (1.0,0.85) rectangle (2.0,1.35);
  \node[font=\tiny, purple!80!black] at (1.5,1.1) {\textbf{叠加区}};
  % 剩余可行域（绿）
  \fill[green!20] (0,2.25) rectangle (4.1,3.1);
  \draw[->, green!60!black, very thick] (0.3,2.67) -- (3.8,2.67);
  \node[font=\tiny, green!60!black] at (2.0,2.95) {仅剩极小可行域};
\end{tikzpicture}

% ============================================================
\begin{tikzpicture}[scale=0.9]
  % 坐标轴
  \draw[->, thick] (0,0) -- (4.3,0) node[right]{\footnotesize $t$/s};
  \draw[->, thick] (0,0) -- (0,3.3) node[above]{\footnotesize $s$/m};
  % 障碍物 1（实心红）
  \fill[red!50] (0.5,0.35) rectangle (1.7,1.05);
  \node[font=\tiny, white] at (1.1,0.7) {\textbf{障碍1}};
  % 障碍物1 安全缓冲（虚线）
  \draw[red!70, dashed, thick] (0.2,0.05) rectangle (2.0,1.35);
  \node[font=\tiny, red!70!black] at (1.7,0.15) {缓冲1};
  % 障碍物 2（实心蓝）
  \fill[blue!50] (1.3,1.15) rectangle (2.5,1.95);
  \node[font=\tiny, white] at (1.9,1.55) {\textbf{障碍2}};
  % 障碍物2 安全缓冲（虚线）
  \draw[blue!70, dashed, thick] (1.0,0.85) rectangle (2.8,2.25);
  \node[font=\tiny, blue!70!black] at (2.5,2.35) {缓冲2};
  % 叠加区高亮（紫）
  \fill[purple!40, opacity=0.75] (1.0,0.85) rectangle (2.0,1.35);
  \node[font=\tiny, purple!80!black] at (1.5,1.1) {\textbf{叠加区}};
  % 剩余可行域（绿）
  \fill[green!20] (0,2.25) rectangle (4.1,3.1);
  \draw[->, green!60!black, very thick] (0.3,2.67) -- (3.8,2.67);
  \node[font=\tiny, green!60!black] at (2.0,2.95) {仅剩极小可行域};
\end{tikzpicture}

% ============================================================
\begin{tikzpicture}[scale=0.9]
  % 坐标轴
  \draw[->, thick] (0,0) -- (4.3,0) node[right]{\footnotesize $t$/s};
  \draw[->, thick] (0,0) -- (0,3.3) node[above]{\footnotesize $s$/m};
  % 障碍物 1（实心红）
  \fill[red!50] (0.5,0.35) rectangle (1.7,1.05);
  \node[font=\tiny, white] at (1.1,0.7) {\textbf{障碍1}};
  % 障碍物1 安全缓冲（虚线）
  \draw[red!70, dashed, thick] (0.2,0.05) rectangle (2.0,1.35);
  \node[font=\tiny, red!70!black] at (1.7,0.15) {缓冲1};
  % 障碍物 2（实心蓝）
  \fill[blue!50] (1.3,1.15) rectangle (2.5,1.95);
  \node[font=\tiny, white] at (1.9,1.55) {\textbf{障碍2}};
  % 障碍物2 安全缓冲（虚线）
  \draw[blue!70, dashed, thick] (1.0,0.85) rectangle (2.8,2.25);
  \node[font=\tiny, blue!70!black] at (2.5,2.35) {缓冲2};
  % 叠加区高亮（紫）
  \fill[purple!40, opacity=0.75] (1.0,0.85) rectangle (2.0,1.35);
  \node[font=\tiny, purple!80!black] at (1.5,1.1) {\textbf{叠加区}};
  % 剩余可行域（绿）
  \fill[green!20] (0,2.25) rectangle (4.1,3.1);
  \draw[->, green!60!black, very thick] (0.3,2.67) -- (3.8,2.67);
  \node[font=\tiny, green!60!black] at (2.0,2.95) {仅剩极小可行域};
\end{tikzpicture}

% ============================================================
\begin{tikzpicture}[scale=0.9]
  % 坐标轴
  \draw[->, thick] (0,0) -- (4.3,0) node[right]{\footnotesize $t$/s};
  \draw[->, thick] (0,0) -- (0,3.3) node[above]{\footnotesize $s$/m};
  % 障碍物 1（实心红）
  \fill[red!50] (0.5,0.35) rectangle (1.7,1.05);
  \node[font=\tiny, white] at (1.1,0.7) {\textbf{障碍1}};
  % 障碍物1 安全缓冲（虚线）
  \draw[red!70, dashed, thick] (0.2,0.05) rectangle (2.0,1.35);
  \node[font=\tiny, red!70!black] at (1.7,0.15) {缓冲1};
  % 障碍物 2（实心蓝）
  \fill[blue!50] (1.3,1.15) rectangle (2.5,1.95);
  \node[font=\tiny, white] at (1.9,1.55) {\textbf{障碍2}};
  % 障碍物2 安全缓冲（虚线）
  \draw[blue!70, dashed, thick] (1.0,0.85) rectangle (2.8,2.25);
  \node[font=\tiny, blue!70!black] at (2.5,2.35) {缓冲2};
  % 叠加区高亮（紫）
  \fill[purple!40, opacity=0.75] (1.0,0.85) rectangle (2.0,1.35);
  \node[font=\tiny, purple!80!black] at (1.5,1.1) {\textbf{叠加区}};
  % 剩余可行域（绿）
  \fill[green!20] (0,2.25) rectangle (4.1,3.1);
  \draw[->, green!60!black, very thick] (0.3,2.67) -- (3.8,2.67);
  \node[font=\tiny, green!60!black] at (2.0,2.95) {仅剩极小可行域};
\end{tikzpicture}
